% !TeX program = xelatex

% Create command `createHandout` only if it doesn't already exists.
% To compile in presentation mode, pass a command createHandout doing nothing.
\providecommand{\createHandout}{handout}

\documentclass[
    aspectratio=169,
    \createHandout
]{beamer}

\usepackage{minted}
\usepackage{xcolor}
\usepackage{tcolorbox}
\usepackage{graphicx}
\usepackage[no-math]{fontspec}
\usepackage{tabularray}
\UseTblrLibrary{diagbox}

% allow to look for themes in another folder
% https://tex.stackexchange.com/a/284157/301699
\makeatletter
\def\beamer@calltheme#1#2#3{%
    \def\beamer@themelist{#2}
    \@for\beamer@themename:=\beamer@themelist\do
    {\usepackage[{#1}]{\beamer@themelocation/#3\beamer@themename}}}

\def\usefolder#1{
    \def\beamer@themelocation{#1}
}
\def\beamer@themelocation{}
\makeatother

\usefolder{..}
\usetheme{cexa-kokkos}

\AtBeginSection{
    \begin{frame}{Outline}
        \renewcommand{\baselinestretch}{0.25}
        \tableofcontents[currentsection]
    \end{frame}
}

\setminted{
    autogobble,
    fontsize=\small,
    bgcolor=lightgray,
    xleftmargin=0.5em,
    xrightmargin=0.5em,
    breaklines,
}

\usemintedstyle[text]{vim}
\setminted[text]{
    bgcolor=lightmain,
    fontsize=\scriptsize,
}

\NewTblrTheme{kokkostable}{
    \SetTblrInner{
        width=\linewidth,
        rowhead=1,
        rows={ht=\baselineskip},
        row{odd}={bg=lightgray},
        row{1}={bg=lightmain},
    }
}

\graphicspath{{../../images/}}

%Information to be included in the title page:
\title{Kokkos basic course}
\author{\scriptsize Mathieu Lobet, Paul Zehner, Paul Gannay, Hariprasad Kannan, Thomas Padioleau}
\institute{CEA}
\date{\today}
\titlegraphic{%
    \includegraphics[height=4em]{kokkos.png}%
    \hspace{1em}%
    \includegraphics[height=4em]{cexa_logo.png}
}
% \titlegraphicbg{\includegraphics[width=\the\paperwidth, trim=0 0 0 0, clip]{custom_bg_image}}

% _____________________________________________________________________________

\begin{document}

\begin{frame}[plain]
    \titlepage
\end{frame}

% _____________________________________________________________________________

\begin{frame}{This course is open source}
    \begin{center}
        \githublink{\url{https://github.com/CExA-project/cexa-kokkos-tutorials}}
    \end{center}
\end{frame}

% _____________________________________________________________________________

\begin{frame}{Prerequisites}
    This course is intended for beginners in Kokkos and performance portability, without strong C++ knowledge

    \vspace{1em}

    \begin{block}{What you still need}
        \begin{itemize}
            \item Basic knowledge of C/C++
            \item Basic knowledge of parallel programming
            \item Basic knowledge of CMake
            \item Basic knowledge of a Linux environment
        \end{itemize}
    \end{block}
\end{frame}

% _____________________________________________________________________________

\begin{frame}{Duration of the course}
    \begin{itemize}
        \item Course + practical work: full day
        \item Course + corrected exercise: half day
        \item Short version: 3 hours
    \end{itemize}
\end{frame}

\begin{frame}{Outline}
    \renewcommand{\baselinestretch}{0.5}
    \tableofcontents
\end{frame}

% _____________________________________________________________________________

\section{Introduction}

% _____________________________________________________________________________

\begin{frame}{Presentation of CExA}
    \begin{columns}
        \begin{column}{0.55\linewidth}
            \begin{center}
                \includegraphics[width=\linewidth]{cexa.png}
            \end{center}
        \end{column}
        \begin{column}{0.45\linewidth}
            \begin{itemize}
                \item \highlight{C}omputing at \highlight{Ex}ascale with \highlight{A}ccelerators at CEA
                \item Alice Recoque exascale supercomputer in 2026, at CEA
                \item Joint effort at CEA to help the migration of HPC codes to GPU
                \begin{itemize}
                    \item Provide GPU programming model (long term, sustainable)
                    \item Support application demonstrators (CEA applications)
                    \item Propose dissemination of information (training, meetings)
                \end{itemize}
            \end{itemize}
        \end{column}
    \end{columns}

\end{frame}

% _____________________________________________________________________________

\begin{frame}{Top500 GPU-CPU system share}
    \begin{center}
        \includegraphics[width=0.8\textwidth]{top500_cpu_gpu_share_history.png}
    \end{center}

    \structure{Note:} Does not include exotic accelerators (Intel Knights Landing, etc.)
\end{frame}

% _____________________________________________________________________________

\begin{frame}{Current supercomputers hardware}
    \begin{center}
        \includegraphics[width=\textwidth]{top10_super_computers.png}

        \caption{Extract of the Top500 (top supercomputers in the world) in November 2025}
    \end{center}
    \begin{itemize}
        \item Heterogeneous nodes: mix of CPU and GPU accelerators
        \item Many CPU (AMD, Intel, NVIDIA) and GPU vendors (AMD, Intel, NVIDIA)
    \end{itemize}
\end{frame}

% _____________________________________________________________________________

\begin{frame}{EuroHPC supercomputers hardware}
    \begin{center}
        \includegraphics[width=1\textwidth]{euroHPC.png}

        \caption{EuroHPC supercomputers}
    \end{center}
\end{frame}

% _____________________________________________________________________________

\begin{frame}{French supercomputers hardware}
    \begin{center}
    \includegraphics[width=0.8\textwidth]{french_super_computers.png}

    \caption{French academic and CEA supercomputers}
    \end{center}
\end{frame}

% _____________________________________________________________________________

\begin{frame}{Basic definition of a supercomputer}
    \begin{center}
        \includegraphics[width=0.8\textwidth]{super-computer_architecture.png}
    \end{center}
    \begin{itemize}
        \item Distributed memory system composed of many compute nodes packed into racks and linked by a high-speed network
        \item Compute nodes are composed of one or more CPUs and one or more accelerators
        \item Most common accelerators today are GPUs (aka GPGPU)
    \end{itemize}
\end{frame}

% _____________________________________________________________________________

\begin{frame}{Zoom on the CPU architecture}
    \begin{columns}
        \begin{column}{0.5\linewidth}
            \begin{center}
                \includegraphics[width=\linewidth]{cpu_architecture.png}
            \end{center}
        \end{column}
        \begin{column}{0.5\linewidth}
            \begin{itemize}
                \item CPUs are for general purpose
                \begin{itemize}
                    \item Sequential task
                    \item Computing
                    \item Operating system
                \end{itemize}
                \item Tens to hundred of cores in biggest processors to perform parallel operations
                \item SIMD (Single Instruction Multiple Data) units to accelerate arithmetic operations
            \end{itemize}
        \end{column}
    \end{columns}
\end{frame}

% _____________________________________________________________________________

\begin{frame}{Zoom on the GPU architecture}
    \begin{columns}
        \begin{column}{0.5\linewidth}
            \begin{center}
                \includegraphics[width=\linewidth]{gpu_architecture.png}
            \end{center}
        \end{column}
        \begin{column}{0.5\linewidth}
            \begin{itemize}
                \item GPUs are specialized
                \begin{itemize}
                    \item Massive parallelism of simple kernels
                    \item Intensive arithmetic computations
                \end{itemize}

                \item Hundreds of computing units, hundred of thousands of threads
                \item Large SIMT vector unit (Single Instruction Multiple Threads) per computing unit
            \end{itemize}
        \end{column}
    \end{columns}
\end{frame}

% _____________________________________________________________________________

\begin{frame}{SIMD versus SIMT}
    \begin{center}
        \includegraphics[width=\textwidth]{SIMD_vs_SIMT.png}
    \end{center}
\end{frame}

% _____________________________________________________________________________

\begin{frame}{Main difference between CPU and GPU}
    \begin{center}
        \includegraphics[width=0.75\textwidth]{cpu_vs_gpu.png}
    \end{center}

    \vspace{-0.5em}

    \begin{columns}
        \begin{column}{0.4\linewidth}
            \begin{itemize}
                \item From about 10 to 200 cores
                \item Complex core architecture and instructions set
                \item General purpose
                \item SIMD model
            \end{itemize}
        \end{column}
        \begin{column}{0.4\linewidth}
            \begin{itemize}
                \item 1,000,000s of threads
                \item Simple threads architecture and instructions set
                \item Specialization
                \item SIMT model
            \end{itemize}
        \end{column}
    \end{columns}
\end{frame}

% _____________________________________________________________________________

\begin{frame}{Host/Device model}
    \begin{columns}[T]
        \begin{column}{0.5\linewidth}
            \begin{center}
                \textbf{CPU}
            \end{center}
            \begin{itemize}
                \item Referred to as the \highlight{Host}
                \item Programs are always started first on it
                \item Orchestrates when the kernels are launched on the GPU and how to make the CPU/GPU memory transfers
            \end{itemize}
        \end{column}
        \begin{column}{0.5\linewidth}
            \begin{center}
                \textbf{GPU}
            \end{center}
            \begin{itemize}
                \item Referred to as the \highlight{Device}
                \item Cannot work standalone
                \item Waits for kernels to execute (offloading)
            \end{itemize}
        \end{column}
    \end{columns}
\end{frame}

\begin{frame}{Host/Device communications}
    \begin{center}
        \includegraphics[width=0.8\linewidth]{cpu_gpu_memory}
    \end{center}
    \begin{itemize}
        \item CPU and GPU are connected with the PCIe bus
        \item The bus memory bandwidth is a fraction of the RAM (CPU memory) or VRAM (GPU memory) bandwidth
        \item Transferring data between host and device is costly
    \end{itemize}
\end{frame}

% _____________________________________________________________________________

\begin{frame}{What's the motivation of putting GPUs in supercomputers?}
    \begin{center}
        \includegraphics[width=0.8\textwidth]{flop_watt_ratio_history_fp64.png}
    \end{center}

    \structure{Question:} How to build the fastest machine with the lowest power consumption and the lowest cost?
\end{frame}

% _____________________________________________________________________________

\begin{frame}{Parallel programming models and libraries in numerical science}
    \begin{center}
        \includegraphics[width=\textwidth]{prog_model.png}
    \end{center}
    \begin{itemize}
        \item Performance portability is one thing
        \item But we will see that developers need more
    \end{itemize}
\end{frame}

% _____________________________________________________________________________

\begin{frame}{Performance, portability and productivity}
    \begin{columns}[T]
        \begin{column}{0.45\linewidth}
            \begin{center}
            \textbf{What \emph{most developers} want}
            \end{center}
            \begin{description}[Performance]
                \item[Performance] Get the best performance on a specific hardware
                \item[Portability] Run the same code on different hardware
                \item[Productivity] Write code quickly and easily, maintain and extend code, leverage architecture optimization
            \end{description}
        \end{column}
        \pause
        \begin{column}{0.55\linewidth}
            \begin{center}
                \textbf{What \emph{developers for science} also want}
            \end{center}
            \begin{description}[Interporability]
                \item[Maturity] Production-ready (i.e. not research product)
                \item[Community] Contributors, support, documentation, examples
                \item[Longevity] Long-term maintenance, bug fixes, updates, sponsorship
                \item[Interporability] Possibility to easily couple with other libraries and tools (IO, Linear Algebra, Machine Learning, etc.)
            \end{description}
        \end{column}
    \end{columns}
\end{frame}

% _____________________________________________________________________________

\begin{frame}{Do \emph{you} need performance portability?}
    \begin{columns}[T]
        \begin{column}{0.6\linewidth}
            \begin{center}
                \textbf{Probably \emph{yes} if you...}
            \end{center}
            \begin{itemize}
                \item Want your code to run on different hardware technologies (CPU, GPU)
                \item Have performance requirements
                \item Cannot afford to maintain and optimize different versions of your code for each possible hardware of the market
                \item Want to focus on algorithmic and not on code development
                \item Want your code to be easily maintainable and readable by others, especially non experts
            \end{itemize}
        \end{column}
        \pause
        \begin{column}{0.4\linewidth}
            \begin{center}
                \textbf{Probably \emph{no} if you...}
            \end{center}
            \begin{itemize}
                \item Need to tune your algorithms to work as fast as possible on a given hardware
                \item Want to explore a new architecture
            \end{itemize}
        \end{column}
    \end{columns}
\end{frame}

\begin{frame}{Programming models (qualitative)}
    \begin{center}
        \includegraphics[width=0.5\textwidth]{p3.png}
    \end{center}
\end{frame}

% _____________________________________________________________________________

\begin{frame}{Kokkos programming model}
    Kokkos is a performance portability parallel programming model built upon the C++-20 standard, designed to abstract already existing parallel programming models

    \vspace{0.5em}

    \begin{center}
        \includegraphics[width=0.9\textwidth]{kokkos_model.png}
    \end{center}
    \begin{itemize}
        \item Kokkos provides data structures and functions based on modern C++ to build your application
        \item Developers never directly use underlying backends (no CUDA kernels, no OpenMP pragmas for instance)
    \end{itemize}
\end{frame}

% _____________________________________________________________________________

\begin{frame}{Kokkos parallelism scope}
    \begin{itemize}
        \item Kokkos only handles node-level parallelism
        \item You still need a model for distributed memory parallelism (e.g.\ MPI, Kokkos-Comm)
    \end{itemize}
\end{frame}

% _____________________________________________________________________________

\begin{frame}{What Kokkos provides you}
    \begin{itemize}
        \item Kokkos' basic capabilities
        \begin{itemize}
            \item Abstraction of most parallel patterns used in HPC (loops, reduction, scan)
            \item Abstracted basic data structures used in science (vectors, multidimensional arrays, sub-arrays, strided arrays, etc)
        \end{itemize}
        \pause
        \item Kokkos' advanced capabilities
        \begin{itemize}
            \item Advanced data structure properties
            \item Thread safety, thread scalability, and atomic operations
            \item Hierarchical patterns for maximizing parallelism
        \end{itemize}
        \pause
        \item Kokkos' libraries
        \begin{itemize}
            \item Profiling, tuning and debugging tools (Kokkos tools)
            \item Python and Fortran interfaces (PyKokkos, Kokkos-Fortran-interop)
            \item Linear algebra (Kokkos-Kernels)
            \item FFT (Kokkos-FFT)
            \item Distributed memory parallelism (Kokkos-Comm)
        \end{itemize}
    \end{itemize}
\end{frame}

% _____________________________________________________________________________

\begin{frame}{Kokkos ecosystem}
    \begin{center}
        \includegraphics[width=0.9\textwidth]{ekokkosystem}
    \end{center}
\end{frame}

% _____________________________________________________________________________

\begin{frame}{Kokkos and the evolution of the C++ standard}
    \begin{center}
        \includegraphics[width=0.9\textwidth]{kokkos-cpp-standard.png}
    \end{center}
\end{frame}

\begin{frame}{Kokkos history and governance}
    \begin{columns}
        \begin{column}{0.45\linewidth}
            \centering

            \includegraphics[width=0.6\linewidth]{hpsf}

            \vspace{2em}

            \includegraphics[width=0.8\linewidth]{kokkos_text}
        \end{column}
        \begin{column}{0.55\linewidth}
            \begin{itemize}
                \item Kokkos was initially developed in American laboratories in 2011
                \item Kokkos is a project of the High Performance Software Foundation (HPSF, part of the Linux Foundation) since 2023
                \item Not a purely American project anymore
                \item Sustainability of the development for the long run
            \end{itemize}
        \end{column}
    \end{columns}
\end{frame}

\begin{frame}{Institutions contributing to Kokkos}
    \centering
    \begin{tabular}{cccc}
        \includegraphics[width=8em, height=3em, keepaspectratio]{anl} &
        \includegraphics[width=8em, height=3em, keepaspectratio]{cea} &
        \includegraphics[width=8em, height=3em, keepaspectratio]{cscs} &
        \includegraphics[width=8em, height=3em, keepaspectratio]{lanl} \\
        \includegraphics[width=8em, height=3em, keepaspectratio]{lbl} &
        \includegraphics[width=8em, height=3em, keepaspectratio]{ornl} &
        \includegraphics[width=8em, height=3em, keepaspectratio]{snl}
    \end{tabular}
\end{frame}

% _____________________________________________________________________________

\begin{frame}{Kokkos useful resources}
    \begin{description}[GitHub repository]
        \item[GitHub repository] \githublink{\url{https://github.com/kokkos}}
        \item[Documentation] \doclink{\url{https://kokkos.org/kokkos-core-wiki}}
        \item[Slack channel] \awesomelink{slack}{\url{https://kokkosteam.slack.com}}
        \item[Cheat sheet] \awesomelink{file}{\url{https://cexa-project.org/documents}}
    \end{description}
\end{frame}

% _____________________________________________________________________________

\section{The basics of Kokkos}

% _____________________________________________________________________________

\begin{frame}{Kokkos}
    \begin{itemize}
        \item Source code in the \githublink{\href{https://github.com/kokkos/kokkos}{Kokkos GitHub repository}}
        \item C++ library, mostly header-only
        \item Requires CMake (no Makefiles!)
        \item Different ways to add it to your project
        \begin{itemize}
            \item External library
            \item Inline build
            \item Spack
        \end{itemize}
    \end{itemize}
    \begin{center}
        \includegraphics[height=3em]{GitHub-logo.png}%
        \hspace{1em}%
        \includegraphics[height=2em]{spack.png}%
    \end{center}
\end{frame}

% _____________________________________________________________________________

\begin{frame}{The notion of backend in Kokkos}
    \begin{itemize}
        \item Kokkos is compiled to target a specific backend
        \item Kokkos is portable until you compile it
        \item You need one compilation per backend
    \end{itemize}
    \begin{center}
        \includegraphics[width=0.6\textwidth]{kokkos_backend.png}
    \end{center}
\end{frame}

% _____________________________________________________________________________

\begin{frame}[fragile]{Compiling Kokkos with the default parameters}
    \begin{columns}
        \begin{column}{0.6\linewidth}
            \begin{itemize}
                \item Defaults to the serial backend
                \item Use the default compiler or specify it with \texttt{-DCMAKE\_CXX\_COMPILER=<command>}
                \item Minimum supported version of the compilers detailed \href{https://kokkos.org/kokkos-core-wiki/get-started/requirements.html}{in the doc}
            \end{itemize}
            \begin{minted}[fontsize=\scriptsize]{bash}
                cmake \
                    -B ${build_dir} \
                    -DCMAKE_BUILD_TYPE=Release \
                    -DCMAKE_CXX_COMPILER=g++
                cmake \
                    --build ${build_dir} \
                    --parallel ${jobs}
                cmake \
                    --install ${build_dir} \
                    --prefix ${path_to_kokkos}
            \end{minted}
        \end{column}
        \pause
        \begin{column}{0.4\linewidth}
            \begin{center}
                \begin{tblr}[theme=kokkostable]{ll}
                    Name       & Cmd.             & Min.\ ver. \\
                    Clang      & \texttt{clang++} & 15.0.0     \\
                    GCC        & \texttt{g++}     & 10.4.0     \\
                    Intel LLVM & \texttt{icpx}    & 2024.2.1   \\
                    NVCC       & \texttt{nvcc}    & 12.2       \\
                    NVHPC      & \texttt{nvc++}   & 22.3       \\
                    ROCM       & \texttt{hipcc}   & 6.2.0
                \end{tblr}
            \end{center}
        \end{column}
    \end{columns}
\end{frame}

% _____________________________________________________________________________

\begin{frame}[fragile]{Compiling Kokkos for a specific backend}
    \begin{columns}
        \begin{column}{0.45\linewidth}
            \begin{itemize}
                \item Add the option \texttt{-DKokkos\_ENABLE\_<BACKEND>=ON} to enable a specific backend
                \item Only one CPU backend and one GPU backend maximum can be compiled at a time (plus Serial backend)
            \end{itemize}
            \begin{minted}{bash}
                cmake \
                    -B ${build_dir} \
                    -DCMAKE_CXX_COMPILER=g++ \
                    -DKokkos_ENABLE_CUDA=ON \
                    ${other_options}
            \end{minted}
        \end{column}
        \pause
        \begin{column}{0.55\linewidth}
            \begin{center}
                \begin{tblr}[theme=kokkostable]{ll}
                    Backend & Option \\
                    Serial & \texttt{-DKokkos\_ENABLE\_SERIAL=ON} \\
                    OpenMP & \texttt{-DKokkos\_ENABLE\_OPENMP=ON} \\
                    Threads & \texttt{-DKokkos\_ENABLE\_THREADS=ON} \\
                    Cuda & \texttt{-DKokkos\_ENABLE\_CUDA=ON} \\
                    HIP & \texttt{-DKokkos\_ENABLE\_HIP=ON} \\
                    SYCL & \texttt{-DKokkos\_ENABLE\_SYCL=ON} \\
                \end{tblr}
            \end{center}
        \end{column}
    \end{columns}
\end{frame}

% _____________________________________________________________________________

\begin{frame}{Warning regarding the backends}
    \begin{itemize}
        \item All backends \highlight{may not have the same level of maturity}
        \item The backend environment (for instance Cuda) must be available on the target machine, \highlight{Kokkos does not install it for you}!
    \end{itemize}
\end{frame}

% _____________________________________________________________________________

\begin{frame}{Compiling Kokkos for a specific architecture}
    \begin{columns}
        \begin{column}{0.4\linewidth}
            \begin{itemize}
                \item Add the architecture option \texttt{-DKokkos\_ARCH\_<NAME>=ON}
                \item CPU architecture is optional, for optimization
                \item GPU architecture is mandatory, deduced if device is present
                \item All CMake options are available \doclink{\href{https://kokkos.org/kokkos-core-wiki/get-started/configuration-guide.html}{in the doc}}
            \end{itemize}
        \end{column}
        \pause
        \begin{column}{0.6\linewidth}
            \begin{center}
                \small
                \begin{tblr}[theme=kokkostable]{ll}
                    GPU model & Option \\
                    AMD MI300A & \texttt{-DKokkos\_ARCH\_AMD\_GFX942\_APU=ON} \\
                    AMD MI300X & \texttt{-DKokkos\_ARCH\_AMD\_GFX942=ON} \\
                    AMD MI250X & \texttt{-DKokkos\_ARCH\_AMD\_GFX90A=ON} \\
                    Intel GPU Max & \texttt{-DKokkos\_ARCH\_INTEL\_PVC=ON} \\
                    NVIDIA H200 & \texttt{-DKokkos\_ARCH\_HOPPER90=ON} \\
                    NVIDIA A100 & \texttt{-DKokkos\_ARCH\_AMPERE80=ON} \\
                    NVIDIA V100 & \texttt{-DKokkos\_ARCH\_VOLTA70=ON} \\
                \end{tblr}
            \end{center}
        \end{column}
    \end{columns}
\end{frame}

% _____________________________________________________________________________

\begin{frame}[fragile]{Real world case: NVIDIA GPU}
    \begin{columns}
        \begin{column}{0.475\linewidth}
            \begin{center}
                \includegraphics[width=\textwidth]{kokkos_a100_compilation.png}
            \end{center}
            \begin{minted}{sh}
                cmake \
                    -B ${build_dir} \
                    -DKokkos_ENABLE_OPENMP=ON \
                    -DKokkos_ARCH_SKX=ON \
                    -DKokkos_ENABLE_CUDA=ON \
                    -DKokkos_ARCH_AMPERE80=ON
            \end{minted}
        \end{column}
        \begin{column}{0.525\linewidth}
            \begin{itemize}
                \item Compiling for a multi-core Xeon Skylake CPU and an NVIDIA A100 GPU
                \item Compiler is GCC (default)
                \item OpenMP option for the multi-core CPU
                \item Skylake option, mostly for vectorization (optional)
                \item Cuda option for the GPU
                \item A100 option, for the correct architecture (mandatory)
            \end{itemize}
        \end{column}
    \end{columns}
\end{frame}

% _____________________________________________________________________________

\begin{frame}[fragile]{Real world case: AMD GPU}
    \begin{columns}
        \begin{column}{0.475\linewidth}
            \begin{center}
                \includegraphics[width=\textwidth]{kokkos_mi300x_compilation.png}
            \end{center}
            \begin{minted}{sh}
                cmake \
                    -B ${build_dir} \
                    -DCMAKE_CXX_COMPILER=hipcc \
                    -DKokkos_ENABLE_OPENMP=ON \
                    -DKokkos_ARCH_SKX=ON \
                    -DKokkos_ENABLE_HIP=ON \
                    -DKokkos_ARCH_AMD_GFX90A=ON
            \end{minted}
        \end{column}
        \begin{column}{0.525\linewidth}
            \begin{itemize}
                \item Compiling for a multi-core Xeon Skylake CPU and an AMD MI300X GPU
                \item Compiler is the HIP compiler
                \item OpenMP option for the multi-core CPU
                \item Skylake option, mostly for vectorization (optional)
                \item HIP option for the GPU
                \item MI300X option, for the correct architecture (mandatory)
            \end{itemize}
        \end{column}
    \end{columns}
\end{frame}

% _____________________________________________________________________________

\begin{frame}[fragile]{Real world case: Intel GPU}
    \begin{columns}
        \begin{column}{0.475\linewidth}
            \begin{center}
                \includegraphics[width=\textwidth]{kokkos_gpu_max_compilation.png}
            \end{center}

            \vspace{-0.5em}

            \begin{minted}[breakafter={=}]{sh}
                cmake \
                    -B ${build_dir} \
                    -DCMAKE_CXX_COMPILER=icpx \
                    -DCMAKE_CXX_FLAGS="-fp-model=precise" \
                    -DKokkos_ENABLE_OPENMP=ON \
                    -DKokkos_ARCH_SKX=ON \
                    -DKokkos_ENABLE_SYCL=ON \
                    -DKokkos_ARCH_INTEL_PVC=ON
            \end{minted}
        \end{column}
        \begin{column}{0.525\linewidth}
            \begin{itemize}
                \item Compiling for a multi-core Xeon Skylake CPU and an Intel GPU Max
                \item Compiler is the Intel LLVM compiler
                \item Add flag to use high precision math operators
                \item OpenMP option for the multi-core CPU
                \item Skylake option, mostly for vectorization (optional)
                \item SYCL option for the GPU
                \item PVC option, for the correct architecture (mandatory)
            \end{itemize}
        \end{column}
    \end{columns}
\end{frame}

% _____________________________________________________________________________

\begin{frame}{How to add Kokkos to your project}
    \begin{itemize}
        \item Kokkos only uses CMake (no Makefiles!)
        \item There are several ways to do it (for lack of a standard C++ dependency manager)
        \begin{itemize}
            \item Build and install it
            \item As a Git submodule
            \item Using CMake FetchContent
            \item With Spack
        \end{itemize}
        \item None of them is the perfect solution
    \end{itemize}
\end{frame}

% _____________________________________________________________________________

\begin{frame}{Different ways to handle the Kokkos dependency}
    \begin{columns}[T]
        \begin{column}{0.25\linewidth}
            \strut Build and install

            \emph{\small\strut(external library)}

            \begin{block}{Pros}
                \begin{itemize}
                    \item You build Kokkos once
                \end{itemize}
            \end{block}
            \begin{alertblock}{Cons}
                \begin{itemize}
                    \item Not flexible
                    \item Not dev-friendly
                    \item Dependency up to the user
                \end{itemize}
            \end{alertblock}
        \end{column}
        \pause
        \begin{column}{0.25\linewidth}
            \strut Git submodule

            \emph{\small\strut(inline build)}

            \begin{block}{Pros}
                \begin{itemize}
                    \item Easy
                    \item Popular
                \end{itemize}
            \end{block}
            \begin{alertblock}{Cons}
                \begin{itemize}
                    \item User may forget to get the submodule
                    \item Not very clean
                \end{itemize}
            \end{alertblock}
        \end{column}
        \pause
        \begin{column}{0.25\linewidth}
            \strut CMake FetchContent

            \emph{\small\strut(inline build)}

            \begin{block}{Pros}
                \begin{itemize}
                    \item Trivial for the user
                \end{itemize}
            \end{block}
            \begin{alertblock}{Cons}
                \begin{itemize}
                    \item Not trivial for the dev to do it right
                    \item Not very clean
                \end{itemize}
            \end{alertblock}
        \end{column}
        \pause
        \begin{column}{0.25\linewidth}
            \strut Spack

            \emph{\small\strut(external library)}

            \begin{block}{Pros}
                \begin{itemize}
                    \item Good for production
                    \item Has dev mode
                    \item Can bring the backend libs
                \end{itemize}
            \end{block}
            \begin{alertblock}{Cons}
                \begin{itemize}
                    \item Not trivial
                \end{itemize}
            \end{alertblock}
        \end{column}
    \end{columns}
\end{frame}

% _____________________________________________________________________________

\section{First steps with Kokkos}

% _____________________________________________________________________________

\begin{frame}[fragile]{Start a Kokkos program}
    \begin{columns}
        \begin{column}{0.6\linewidth}
            \begin{minted}{C++}
                #include <Kokkos_Core.hpp>

                int main(int argc, char* argv[]) {
                    Kokkos::initialize(argc, argv);
                    {
                        // Your code here
                    }
                    Kokkos::finalize();
                }
            \end{minted}
        \end{column}
        \begin{column}{0.4\linewidth}
            \begin{itemize}
                \item Include the Kokkos header file (like any other library)
                \item Use the Kokkos namespace \texttt{Kokkos::}
                \item Initialize and finalize Kokkos (like MPI)
                \item Takes \texttt{argc} and \texttt{argv} to be interacted from the command line
            \end{itemize}
        \end{column}
    \end{columns}
\end{frame}

% _____________________________________________________________________________

\begin{frame}[fragile]{Start a Kokkos program (RAII)}
    \begin{columns}
        \begin{column}{0.6\linewidth}
            \begin{minted}{C++}
                #include <Kokkos_Core.hpp>

                int main(int argc, char* argv[]) {
                    Kokkos::ScopeGuard guard(argc, argv);

                    // Your code here
                }
            \end{minted}
        \end{column}
        \begin{column}{0.4\linewidth}
            \begin{itemize}
                \item \texttt{ScopeGuard} is a class that ensures that \texttt{Kokkos::initialize} and \texttt{Kokkos::finalize} are called correctly
                \item Resource Acquisition Is Initialization (RAII) pattern
            \end{itemize}
        \end{column}
    \end{columns}
\end{frame}

% _____________________________________________________________________________

\begin{frame}[fragile]{Use CMake to compile your Kokkos program (simple approach)}
    \begin{columns}
        \begin{column}{0.6\linewidth}
            \begin{minted}{cmake}
                cmake_minimum_required(VERSION 3.22)

                project(my_kokkos_project LANGUAGES CXX)

                find_package(Kokkos REQUIRED)

                add_executable(proj main.cpp)
                target_link_libraries(proj Kokkos::kokkos)
            \end{minted}
        \end{column}
        \begin{column}{0.4\linewidth}
            \begin{itemize}
                \item \texttt{find\_package} command to find Kokkos
                \item \texttt{REQUIRED} aborts the configuration if Kokkos is not found
                \item \texttt{target\_link\_libraries} command to link your program with Kokkos
                \item This is the "Kokkos is built and installed" and "Spack" approaches
            \end{itemize}
        \end{column}
    \end{columns}
\end{frame}

% _____________________________________________________________________________

\begin{frame}[fragile]{Making the \texttt{find\_package} part more flexible}
    \begin{columns}
        \begin{column}{0.6\linewidth}
            \begin{minted}[breakafter=/]{cmake}
                find_package(Kokkos QUIET)
                if(NOT Kokkos_FOUND)
                    include(FetchContent)
                    FetchContent_Declare(
                        kokkos
                        URL https://github.com/kokkos/kokkos/releases/download/5.2.2/kokkos-5.2.2.zip
                    )
                    FetchContent_MakeAvailable(kokkos)
                endif()
            \end{minted}
        \end{column}
        \begin{column}{0.4\linewidth}
            \begin{itemize}
                \item \texttt{find\_package} command to find Kokkos if it's already installed
                \item \texttt{QUIET} disables informational messages, even if the package is not found
                \item \texttt{FetchContent\_*} commands to download Kokkos if needed
                \item This is the "Kokkos is managed using CMake" approach
            \end{itemize}
        \end{column}
    \end{columns}
\end{frame}

% _____________________________________________________________________________

\begin{frame}[fragile]{Build your Kokkos program}
    \begin{columns}
        \begin{column}{0.6\linewidth}
            \begin{minted}{bash}
                # configure
                cmake \
                    -B ${build_dir} \
                    -DCMAKE_CXX_COMPILER=${compiler} \
                    -DKokkos_ROOT=${path_to_kokkos} \
                    ${other_options}

                # build
                cmake \
                    --build ${build_dir} \
                    --parallel ${jobs}
            \end{minted}
        \end{column}
        \begin{column}{0.4\linewidth}
            \begin{itemize}
                \item If Kokkos is installed, it propagates its compilation flags to your program
                \item \texttt{-DKokkos\_ROOT} to tell where Kokkos is installed
                \item Otherwise, you specify Kokkos options along with the ones of your program
                \item You may want to add \texttt{-DCMAKE\_BUILD\_TYPE=<Mode>}
                \item \texttt{--parallel N} to build in parallel
            \end{itemize}
        \end{column}
    \end{columns}
\end{frame}

% _____________________________________________________________________________

\begin{exerciseframe}{Exercise 1: First Kokkos program}
    \begin{columns}
        \begin{column}{0.5\linewidth}
            \begin{center}
                \includegraphics[width=0.9\textwidth]{sleeping_otter.png}
            \end{center}

            Go to the exercise \githublink{\href{https://github.com/CExA-project/cexa-kokkos-tutorials/tree/main/exercises/01_first_program}{01\_first\_program}} and follow the instructions
        \end{column}
        \begin{column}{0.5\linewidth}
            \begin{block}{Goal of this exercise}
                \begin{itemize}
                    \item Write a simple Kokkos program
                    \item Compile the program with CMake for different backends
                    \item Get familiar with the compilation process
                \end{itemize}
            \end{block}
        \end{column}
    \end{columns}
\end{exerciseframe}

% _____________________________________________________________________________

\section[Data container]{Data container}

% _____________________________________________________________________________

\begin{frame}[fragile]{Kokkos' abstracted data container}
    \begin{itemize}
        \item Why is it important to have an abstracted data container?
        \begin{itemize}
            \item<1-> No need to allocate or deallocate memory by hand
            \item<2-> Unified semantic and portable memory management (CPU and GPU)
            \item<3-> Vendor-specific memory allocation is hidden
            \item<4-> Advanced capabilities (abstracted layout, subarray, multidimensionality, etc.)
            \item<5-> Safety, through runtime or compile time checks
        \end{itemize}
        \item<6-> Kokkos provides the View class
        \begin{itemize}
            \item<7-> View is an abstraction of the notion of container and multidimensional array
            \item<8-> Brings portable Python numpy/Fortran-like syntax
        \end{itemize}
    \end{itemize}
\end{frame}

% _____________________________________________________________________________

\begin{frame}[fragile]{Creating a View}
    \begin{columns}
        \begin{column}{0.6\linewidth}
            \begin{minted}{C++}
                Kokkos::View<float*>
                    view1("dynamic size", 10);

                Kokkos::View<int[100]>
                    view2("static size");

                Kokkos::View<double*[10]>
                    view3("static and dynamic size", 10);

                Kokkos::View<double***[10]>
                    view4("four dimensions", 10, 10, 10);
            \end{minted}
        \end{column}
        \begin{column}{0.4\linewidth}
            \begin{itemize}
                \item View is a template class
                \item Contains any C++ type (\texttt{int}, \texttt{float}, \texttt{double}, etc)
                \item Number of dimensions, fixed at compile time (no more than 7)
                \item Size of each dimension determined at...
                \begin{description}[\texttt{[n]}]
                    \item[\texttt{*}] runtime
                    \item[\texttt{[n]}] compile time
                \end{description}
            \end{itemize}
        \end{column}
    \end{columns}
\end{frame}

% _____________________________________________________________________________

\begin{frame}[fragile]{Creating a View (2)}
    \begin{columns}
        \begin{column}{0.6\linewidth}
            \begin{minted}{C++}
                Kokkos::View<float*>
                    view1("dynamic size", 10);

                Kokkos::View<int[100]>
                    view2("static size");

                Kokkos::View<double*[10]>
                    view3("static and dynamic size", 10);

                Kokkos::View<double***[10]>
                    view4("four dimensions", 10, 10, 10);
            \end{minted}
        \end{column}
        \begin{column}{0.4\linewidth}
            \begin{itemize}
                \item Name (useful for debugging!)
                \item Check at compile time and runtime (out of bound, same size for operations, ...)
            \end{itemize}
        \end{column}
    \end{columns}
\end{frame}

% _____________________________________________________________________________

\begin{frame}[fragile]{Accessing the data}
    \begin{columns}
        \begin{column}{0.6\linewidth}
            \begin{minted}{C++}
                const int Nx = 1000;
                const int Ny = 2000;

                Kokkos::View<int**>
                    my_matrix("matrix", Nx, Ny);

                for (int i = 0; i < Nx; i++) {
                    for (int j = 0; j < Ny; j++) {
                        my_matrix(i, j) = i + j;
                    }
                }
            \end{minted}
        \end{column}
        \begin{column}{0.4\linewidth}
            \begin{itemize}
                \item Data are accessed using the parenthesis operator \texttt{(i, j, k, ...)}
                \item Raw data can be accessed using the \texttt{data()} method (not recommended)
            \end{itemize}
        \end{column}
    \end{columns}
\end{frame}

% _____________________________________________________________________________

\begin{frame}[fragile]{Basic methods to manage views}
    \begin{columns}
        \begin{column}{0.45\linewidth}
            \begin{minted}{C++}
                Kokkos::View<int ***>
                    view("view", 10, 20, 30);

                int rank = view.rank();
                assert(rank == 3);

                int size_x = view.extent(0);
                assert(size_x == 10);

                int size = view.size();
                assert(size == 6000);
            \end{minted}
        \end{column}
        \begin{column}{0.55\linewidth}
            \begin{description}[\texttt{extent(dim)}]
                \item[\texttt{rank()}] Return the rank of the view (number of dimensions)
                \item[\texttt{extent(dim)}] Return the number of elements of a specific dimension
                \item[\texttt{size()}] Return the total number of elements
            \end{description}

            \vspace{0.5em}

            And many others not covered in this talk
        \end{column}
    \end{columns}
\end{frame}

% _____________________________________________________________________________

\begin{frame}{Where are a View data stored?}
    \begin{center}
        \includegraphics[width=0.65\textwidth]{view_memory.png}
    \end{center}
    \pause
    \begin{itemize}
        \item A View data is allocated only in one specific memory (Host or Device), never both
        \pause
        \item Suppose we created a view on the Host
    \end{itemize}
\end{frame}

% _____________________________________________________________________________

\begin{frame}{Where are a View data stored? CPU version}
    \begin{center}
        \includegraphics[width=0.65\textwidth]{host_view_memory.png}
    \end{center}
    \begin{itemize}
        \item If Kokkos is compiled with a \highlight{CPU backend only}, the View data is allocated in the \highlight{Host memory} by default
        \item View comes with metadata (number of dimensions, extents, etc.) and data
        \item Host view data usually \highlight{cannot} be accessed from the Device
    \end{itemize}
\end{frame}

% _____________________________________________________________________________

\begin{frame}{Where are a View data stored? GPU version}
    \begin{center}
        \includegraphics[width=0.65\textwidth]{device_view_memory.png}
    \end{center}

    \vspace{-0.5em}

    \begin{itemize}
        \item If Kokkos is compiled with a \highlight{GPU backend}, the View data is allocated in the \highlight{Device memory} by default
        \item Metadata are on the Host
        \item Device view data usually \highlight{cannot} be accessed on the Host (we'll see later how to transfer data)
    \end{itemize}
\end{frame}

% _____________________________________________________________________________

\begin{frame}{The concept of memory space}
    \begin{center}
        \includegraphics[width=0.7\linewidth]{memory_space.png}
    \end{center}
    \begin{itemize}
        \item Kokkos abstraction for where the data are allocated: \highlight{memory space}
        \item A View is always associated, at compile time, to a specific memory space (Host or Device)
        \item By default, view data is allocated
        \begin{itemize}
            \item On the Host if no GPU backend is available
            \item On the Device otherwise
        \end{itemize}
    \end{itemize}
\end{frame}

% _____________________________________________________________________________

\begin{frame}[fragile]{How to set a memory space}
    \begin{columns}
        \begin{column}{0.5\linewidth}
            \begin{minted}{C++}
                Kokkos::View<int**, MemorySpace>
                    my_matrix("matrix", 10, 10);
            \end{minted}
        \end{column}
        \begin{column}{0.5\linewidth}
            \begin{itemize}
                \item Template argument \texttt{MemorySpace} to specify the memory space when creating a view
                \item Several ones available
                \item Using a backend-specific memory space is possible, but it breaks portability
                \item Common values are usually sufficient for basic usages
            \end{itemize}
        \end{column}
    \end{columns}
\end{frame}

% _____________________________________________________________________________

\begin{frame}{Notion of View layout}
    \begin{center}
        \includegraphics[width=0.8\textwidth]{layout_right_left.png}
    \end{center}
    \begin{itemize}
        \item Layout is the way multidimensional indices map to linear memory
        \item Kokkos provides an abstraction of the data layout
        \item The default layout of a View depends on the backend (Host or Device)
    \end{itemize}
\end{frame}

% _____________________________________________________________________________

\begin{frame}[fragile]{How to set a View layout}
    \begin{columns}
        \begin{column}{0.5\linewidth}
            \begin{minted}{C++}
                Kokkos::View<int**, Layout>
                    my_matrix("matrix", 10, 10);
            \end{minted}
        \end{column}
        \begin{column}{0.5\linewidth}
            \begin{itemize}
                \item Template argument \texttt{Layout} to specify the layout when creating a view
                \item Several ones available (Left, Right, Stride, ...)
                \item Possible to define your own custom layout
                \item Default values are usually sufficient for basic usages
            \end{itemize}
        \end{column}
    \end{columns}
\end{frame}

% _____________________________________________________________________________

\begin{frame}{More about views}
    \begin{itemize}
        \item Allocations only happen when explicitly specified
        \pause
        \item Reference counting is used for automatic deallocation (like shared pointers)
        \pause
        \item Copy construction and assignment are shallow: you pass Views by value, not by pointer (e.g. C arrays); you can still pass by constant reference
    \end{itemize}
\end{frame}

% _____________________________________________________________________________

\begin{frame}[fragile]{Copy and assignation}
    \begin{columns}
        \begin{column}{0.6\linewidth}
            \begin{minted}{C++}
                Kokkos::View<int*[5]> a("a", 10), b("b", 15);
                a = b;
                Kokkos::View<int**> c(b);

                a(0, 2) = 1;
                b(0, 2) = 2;
                c(0, 2) = 3;

                std::cout << a(0, 2) << std::endl;
            \end{minted}
        \end{column}
        \begin{column}{0.4\linewidth}
            \structure{Question:} What gets printed?

            \pause

            \structure{Answer:} 3
        \end{column}
    \end{columns}
\end{frame}

% _____________________________________________________________________________

\begin{exerciseframe}{Exercise 2: How to use and manage Kokkos Views}
    \begin{columns}
        \begin{column}{0.5\linewidth}
            \begin{center}
                \includegraphics[width=0.9\textwidth]{sleeping_otter.png}
            \end{center}

            Go to the exercise \githublink{\href{https://github.com/CExA-project/cexa-kokkos-tutorials/tree/main/exercises/02_basic_view}{02\_basic\_view}} and follow the instructions
        \end{column}
        \begin{column}{0.5\linewidth}
            \begin{block}{Goal of this exercise}
                \begin{itemize}
                    \item Create a View
                    \item Manage the view to inspect its metadata
                \end{itemize}
            \end{block}
        \end{column}
    \end{columns}
\end{exerciseframe}

% _____________________________________________________________________________

\begin{frame}{How to access data living on the Device from the Host?}
    \begin{center}
        \includegraphics<1|handout:0>[width=0.9\textwidth]{device_memory_access_1.png}
        \includegraphics<2>[width=0.9\textwidth]{device_memory_access_2.png}
    \end{center}

    \onslide<2>{\structure{Problem:} You cannot!}
\end{frame}

% _____________________________________________________________________________

\begin{frame}[fragile]{All you need is a mirror}
    \begin{columns}
        \begin{column}{0.6\linewidth}
            \begin{minted}{C++}
                Kokkos::View<int**>
                    device_matrix("matrix", 10, 10);

                auto host_matrix =
                    Kokkos::create_mirror(device_matrix);
            \end{minted}

        \end{column}
        \begin{column}{0.4\linewidth}
            \begin{itemize}
                \item A \highlight{mirror} is a view in a different memory space
                \item Similar to its original view for the rest (shape, layout, etc.)
                \item Used to create a mirror of a Device view on the Host
                \item If the original view is already on the Host (CPU backend), it creates a mirror without allocating memory (shallow copy of the original view)
                \item \texttt{create\_mirror} allocates the memory but it \highlight{doesn't copy the data}!
            \end{itemize}
        \end{column}
    \end{columns}
\end{frame}

% _____________________________________________________________________________

\begin{frame}{How to access data living on the Device from the Host?}
    \begin{center}
        \includegraphics[width=0.9\textwidth]{device_mirror_view.png}
    \end{center}

    \structure{Solution:} With a mirror view!
\end{frame}

% _____________________________________________________________________________

\begin{frame}{How to access data living on the Host from the Host?}
    \begin{center}
        \includegraphics[width=0.9\textwidth]{host_mirror_view.png}
    \end{center}
    \begin{itemize}
        \item For portability reason, we always use a mirror view whether we compile with a CPU or GPU backend
        \item With a CPU backend only, the mirror view is allocated in the same memory space as the original view
    \end{itemize}

    \structure{Problem:} Memory is duplicated unnecessarily!
\end{frame}

% _____________________________________________________________________________

\begin{frame}[fragile]{All you need is a (different kind of) mirror}
    \begin{columns}
        \begin{column}{0.525\linewidth}
            \begin{minted}{C++}
                Kokkos::View<int**>
                    device_matrix("matrix", 10, 10);

                auto host_matrix =
                    Kokkos::create_mirror_view(
                        device_matrix
                    );
            \end{minted}
        \end{column}
        \begin{column}{0.475\linewidth}
            \begin{itemize}
                \item \texttt{create\_mirror\_view} creates a mirror without allocating memory if the mirror view is allocated in the same memory space as the original view
                \item In that case, the mirror view is a shallow copy of the original view (i.e. it targets the same data)
            \end{itemize}
        \end{column}
    \end{columns}
\end{frame}

% _____________________________________________________________________________

\begin{frame}[fragile]{Mirror functions wrap up}
    \begin{minted}{C++}
        auto mirror = Kokkos::create_mirror(view);
        auto mirror_view = Kokkos::create_mirror_view(view);
    \end{minted}
    \begin{center}
        \begin{tblr}[theme=kokkostable]{colspec=lll, row{2}={bg=lightmain}}
            & \SetCell[c=2]{l} Kokkos compiled on \\
            Mirror function & CPU backend & GPU backend \\
            \texttt{Kokkos::create\_mirror} & Allocate & Allocate \\
            \texttt{Kokkos::create\_mirror\_view} & Don't allocate & Allocate \\
        \end{tblr}
    \end{center}
\end{frame}

% _____________________________________________________________________________

\begin{frame}{How to access data living on the Host from the Host?}
    \begin{center}
        \includegraphics[width=0.9\textwidth]{host_create_mirror_view.png}
    \end{center}

    \structure{Solution:} Use \texttt{create\_mirror\_view}!
\end{frame}

% _____________________________________________________________________________

\begin{frame}[fragile]{Deep copy data between memory spaces}
    \begin{columns}
        \begin{column}{0.5\linewidth}
            \begin{minted}{C++}
                Kokkos::View<int**>
                    device_view("view", Nx, Ny);

                auto host_view =
                    Kokkos::create_mirror_view(
                        device_view
                    );

                initialize(host_view);

                // from Host to Device
                Kokkos::deep_copy(
                    device_view,
                    host_view
                );
            \end{minted}
        \end{column}
        \begin{column}{0.5\linewidth}
            \begin{itemize}
                \item \texttt{deep\_copy(destination, source)} copy \texttt{source} to \texttt{destination}, even if the data are in different memory spaces
                \item Beware of \highlight{arguments order!}
                \item Views must be similar (same shape, layout, etc.)
                \item If \texttt{deep\_copy} detects that the destination view is a shallow copy of the source view, it does nothing
                \item \texttt{deep\_copy(destination, source)} does an automatic fence
            \end{itemize}
        \end{column}
    \end{columns}
\end{frame}

% _____________________________________________________________________________

\begin{frame}{How to deep copy between the Host and the Device?}
    \begin{center}
        \includegraphics[width=0.9\textwidth]{host_device_deep_copy.png}
    \end{center}

    \structure{Solution:} With \texttt{create\_mirror\_view} and \texttt{deep\_copy}!
\end{frame}

% _____________________________________________________________________________

\begin{exerciseframe}{Exercise 3: Mirror Views and deep copy}
    \begin{columns}
        \begin{column}{0.5\linewidth}
            \begin{center}
                \includegraphics[width=0.9\textwidth]{sleeping_otter.png}
            \end{center}

            Go to the exercise \githublink{\href{https://github.com/CExA-project/cexa-kokkos-tutorials/tree/main/exercises/03_deep_copy}{03\_deep\_copy}} and follow the instructions
        \end{column}
        \begin{column}{0.5\linewidth}
            \begin{block}{Goal of this exercise}
                \begin{itemize}
                    \item Create and manage a mirror view
                    \item Use the \texttt{deep\_copy} function
                    \item Observe the differences between the two mirror functions
                    \begin{itemize}
                        \item Timing
                        \item Memory (on Unixes only)
                    \end{itemize}
                    \item Increase the size of the problem to make the differences more visible
                \end{itemize}
            \end{block}
        \end{column}
    \end{columns}
\end{exerciseframe}

% _____________________________________________________________________________

\begin{frame}{Subviews description}
    \begin{itemize}
        \item A subview is a ’slice’ of a View that behaves as a View
        \pause
        \begin{itemize}
            \item Same syntax as a View - access data using (multi-)index entries
            \pause
            \item The ’slice’ and original View point to the same data - no extra
            memory allocation or copying
        \end{itemize}
        \pause
        \item Can be constructed on host or within a kernel (no allocation
        of memory occurs)
        \pause
        \item Similar capability as provided by Matlab, Fortran, Python, etc.
        using ’colon’ notation
    \end{itemize}
\end{frame}

% _____________________________________________________________________________

\begin{frame}[fragile]{Usage demo}
    Begin with a View:
    \begin{minted}{C++}
        Kokkos::View<double ***> v("v", N0, N1, N2);
    \end{minted}
    \pause
    Say we want a 2-dimensional slice at an index i0 in the first
    dimension - that is, in Matlab/Fortran/Python notation:
    \begin{minted}{python}
        slicei0 = v[i0, :, :]
    \end{minted}
    \pause
    This can be accomplished in Kokkos using a subview as follows:
    \begin{minted}{C++}
        auto slice = Kokkos::subview(v, i0, Kokkos::ALL, Kokkos::ALL);
    \end{minted}
\end{frame}

% _____________________________________________________________________________

\begin{frame}[fragile]{Syntax}
    \begin{minted}{C++}
        Kokkos::subview(view, arg0, ...);
    \end{minted}
    \begin{itemize}
        \item \texttt{view}: First argument to the subview is the view of which a slice will be taken
        \item \texttt{argN}: Slice info for rank $N$ - provide same number of arguments as rank
        \item Options for \texttt{argN}:
        \begin{description}[partial range]
            \item[index] integral type single value
            \item[partial range] \texttt{std::pair} or \texttt{Kokkos::pair} of integral types to provide sub-range of a rank’s range $[0, N)$
            \item[full range] use \texttt{Kokkos::ALL} rather than providing the full range as a pair
        \end{description}
    \end{itemize}
\end{frame}

% _____________________________________________________________________________

\begin{frame}[fragile]{Subviews suggested usage}
    \begin{itemize}
        \item Use \texttt{auto} to determine the return type of a subview
        \item A subview can help with encapsulation - e.g. can pass into functions expecting a lower-dimensional View
        \item Use \texttt{Kokkos::pair} for partial ranges if subview created within a kernel
        \item Avoid usage if very few data accesses will be made to the subview
        \begin{itemize}
            \item Construction of subview costs 20-40 operations
        \end{itemize}
    \end{itemize}
\end{frame}

% _____________________________________________________________________________

\section{Parallel loops}

% _____________________________________________________________________________

\begin{frame}[fragile]{Anatomy of a loop}
    \begin{columns}
        \begin{column}{0.5\linewidth}
            \begin{minted}[texcomments]{C++}
                for (int i = 0; i < N; i++) {
                    A(i) = B(i) + C(i) * D(i);
                }
            \end{minted}
        \end{column}
        \begin{column}{0.5\linewidth}
            \begin{itemize}
                \item \mintinline{C++}{for} is the loop pattern (for, reduction, scan, graph)
                \item \mintinline{C++}{int i = 0; i < N; i++} is the execution policy
                \item \mintinline{C++}{A(i) = B(i) + C(i) * D(i)} is the kernel
                \begin{itemize}
                    \item No iteration order dependencies
                    \item No side effects
                \end{itemize}
            \end{itemize}
        \end{column}
    \end{columns}
\end{frame}

% _____________________________________________________________________________

\begin{frame}[fragile]{Comparison of Kokkos and OpenMP parallel loops}
    \begin{columns}[T]
        \begin{column}{0.5\linewidth}
            \begin{center}
                \textbf{OpenMP version}
            \end{center}
            \begin{minted}{C++}
                #pragma omp parallel for
                for (int i = 0; i < N; i++) {
                    A(i) = B(i) + C(i) * D(i);
                }
            \end{minted}
            \begin{block}{Note}
                Only works on CPUs (needs the \texttt{target} directive for GPUs)
            \end{block}
        \end{column}
        \pause
        \begin{column}{0.5\linewidth}
            \begin{center}
                \textbf{Kokkos version}
            \end{center}
            \begin{minted}{C++}
                Kokkos::parallel_for(
                    "my_loop",
                    N,
                    KOKKOS_LAMBDA (int i) {
                        A(i) = B(i) + C(i) * D(i);
                    }
                );
            \end{minted}

            \vspace{-0.75em}

            \begin{block}{Note}
                Works on CPUs and GPUs depending on the compiled backend
            \end{block}
        \end{column}
    \end{columns}
\end{frame}

% _____________________________________________________________________________

\begin{frame}[fragile]{Anatomy of a Kokkos parallel loop}
    \begin{columns}
        \begin{column}{0.5\linewidth}
            \begin{minted}{C++}
                Kokkos::parallel_for(
                    "my_loop",
                    N,
                    KOKKOS_LAMBDA (int i) {
                        A(i) = B(i) + C(i) * D(i);
                    }
                );
            \end{minted}
        \end{column}
        \begin{column}{0.5\linewidth}
            \begin{itemize}
                \item \mintinline{C++}{parallel_for} is the loop pattern
                \item \mintinline{C++}{"my_loop"} is the name of the loop (bonus for debugging!)
                \item \mintinline{C++}{N} and \mintinline{C++}{int i} is the execution policy (can't be simpler)
                \item \mintinline[breaklines]{C++}{KOKKOS_LAMBDA (int i) {/*...*/}} is the kernel
            \end{itemize}
        \end{column}
    \end{columns}

    \pause

    \vspace{1em}
    \structure{Question:} What is this \texttt{KOKKOS\_LAMBDA} thingy?
\end{frame}

% _____________________________________________________________________________

\begin{frame}[fragile]{Notion of C++ lambdas}
    \begin{columns}
        \begin{column}{0.5\linewidth}
            \begin{minted}{C++}
                auto kernel =
                    KOKKOS_LAMBDA (int i) {
                        A(i) = B(i) + C(i) * D(i);
                    };
            \end{minted}
        \end{column}
        \begin{column}{0.5\linewidth}
            \begin{itemize}
                \item A lambda is a C++ anonymous function (i.e. a function as an object)
                \item It has explicit arguments (\mintinline{C++}|int i|) and captured objects (\mintinline{C++}|A|, \mintinline{C++}|B|, \mintinline{C++}|C|, \mintinline{C++}|D|)
                \item \texttt{KOKKOS\_LAMBDA} is a macro that creates a lambda function with the correct capture for you (and a bit more)
                \item For beginners, no need to understand C++ lambdas to use Kokkos (the intermediate course gives more details)
            \end{itemize}
        \end{column}
    \end{columns}
\end{frame}

% _____________________________________________________________________________

\begin{frame}{Now, where is my parallel loop executed?}
    \begin{center}
        \includegraphics[width=0.8\linewidth]{execution_space.png}
    \end{center}
    \begin{itemize}
        \item Kokkos abstraction for where the loop is executed : \highlight{execution space}
        \pause
        \item Parallel code (i.e. inside a \texttt{parallel\_*}) is by default executed
        \begin{itemize}
            \item By the Host if no GPU backend is available
            \item By the Device otherwise
        \end{itemize}
        \pause
        \item Serial code (the rest) is always executed on the Host
    \end{itemize}
\end{frame}

% _____________________________________________________________________________

\begin{frame}[fragile]{Asynchronous execution on GPU}
    \begin{columns}
        \begin{column}{0.5\linewidth}
            \begin{minted}{C++}
                // parallel loop on Device
                Kokkos::parallel_for(
                    "my_loop",
                    N,
                    KOKKOS_LAMBDA (int i) {
                        data(i) = // ...
                    }
                );

                // Host code
                host_function(data);
            \end{minted}
        \end{column}
        \begin{column}{0.5\linewidth}
            \begin{itemize}
                \item With a GPU backend, the parallel loop is executed asynchronously on the GPU
                \item The loop is launched and the program continues
                \item Asynchronism is a complex concept for beginners, not covered in this talk (but in the intermediate one)
            \end{itemize}
        \end{column}
    \end{columns}

    \pause

    \vspace{1em}
    \structure{Problem:} What if I need the result of the parallel loop in the host code?
\end{frame}

% _____________________________________________________________________________

\begin{frame}[fragile]{All you need is a fence}
    \begin{columns}
        \begin{column}{0.5\linewidth}
            \begin{minted}{C++}
                // parallel loop on Device
                Kokkos::parallel_for(
                    "my_loop",
                    N,
                    KOKKOS_LAMBDA (int i) {
                        data(i) = // ...
                    }
                );

                Kokkos::fence("my_loop fence");

                // Host code
                host_function(data);
            \end{minted}
        \end{column}
        \begin{column}{0.5\linewidth}
            \begin{itemize}
                \item \texttt{fence} to wait for asynchronous work to be completed
                \item Equivalent to OpenMP \texttt{barrier} or to \texttt{MPI\_Wait}
                \item Fence label (bonus for debugging!)
            \end{itemize}
        \end{column}
    \end{columns}

    \vspace{1em}
    \structure{Solution:} With a fence
\end{frame}

% _____________________________________________________________________________

\begin{exerciseframe}{Exercise 4: My first parallel loop}
    \begin{columns}
        \begin{column}{0.5\linewidth}
            \begin{center}
                \includegraphics[width=0.9\textwidth]{sleeping_otter.png}
            \end{center}

            Go to the exercise \githublink{\href{https://github.com/CExA-project/cexa-kokkos-tutorials/tree/main/exercises/04_parallel_loop}{04\_parallel\_loop}} and follow the instructions
        \end{column}
        \begin{column}{0.5\linewidth}
            \begin{block}{Goal of this exercise}
                \begin{itemize}
                    \item Write a simple parallel loop
                    \item Measure the difference of performance between CPU and GPU
                    \item Increase the size of the problem to make the difference more visible
                    \item At which size the GPU is more efficient?
                \end{itemize}
            \end{block}
        \end{column}
    \end{columns}
\end{exerciseframe}

% _____________________________________________________________________________

\section{Extending loop policies}

% _____________________________________________________________________________

\begin{frame}[fragile]{How to extend loop policies?}
    \begin{columns}
        \begin{column}{0.55\linewidth}
            \begin{minted}{C++}
                Kokkos::RangePolicy<ExecutionSpace>
                    policy(start_index, end_index);
            \end{minted}
        \end{column}
        \begin{column}{0.45\linewidth}
            \begin{itemize}
                \item Kokkos way to tell how a loop is executed: \highlight{execution policy}
                \item \mintinline{c++}|RangePolicy| for single-dimensional loops
                \item \mintinline{c++}|ExecutionSpace| for specifying where the loop is executed (defaults to \mintinline{c++}|DefaultExecutionSpace|)
                \item \mintinline{c++}|start_index| and \mintinline{c++}|end_index| are the beginning (inclusive) and the end (exclusive) of the loop
            \end{itemize}
        \end{column}
    \end{columns}
\end{frame}

% _____________________________________________________________________________

\begin{frame}[fragile]{Loop policy shortcut}
    \begin{columns}[T]
        \begin{column}{0.45\linewidth}
            \begin{center}
                \textbf{Shortcut version\strut}
            \end{center}

            \begin{minted}{C++}
                Kokkos::parallel_for(
                    "my_loop",
                    N,
                    KOKKOS_LAMBDA (int i) {
                        /* ... */
                    }
                );
            \end{minted}
        \end{column}
        \pause
        \begin{column}{0.55\linewidth}
            \begin{center}
                \textbf{Explicit version (same)\strut}
            \end{center}

            \begin{minted}{C++}
                Kokkos::parallel_for(
                    "my_loop",
                    Kokkos::RangePolicy(0, N),
                    KOKKOS_LAMBDA (int i) {
                        /* ... */
                    }
                );
            \end{minted}
        \end{column}
    \end{columns}
\end{frame}

% _____________________________________________________________________________

\begin{frame}{Execution space}
    \begin{itemize}
        \item Execution space to control where the loop is executed (CPU or GPU)
        \item Several ones available
        \item \texttt{Kokkos::DefaultExecutionSpace} and \texttt{Kokkos::DefaultHostExecutionSpace} are enough for beginners
        \item Using a backend-specific execution space is possible, but it breaks portability
    \end{itemize}
    \pause
    \begin{center}
        \begin{tblr}[theme=kokkostable]{colspec=lll, row{2}={bg=lightmain}}
            & \SetCell[c=2]{l} Kokkos compiled on \\
            Execution space & CPU backend & GPU backend \\
            \texttt{Kokkos::DefaultExecutionSpace} & On CPU & On GPU \\
            \texttt{Kokkos::DefaultHostExecutionSpace} & On CPU & On CPU \\
        \end{tblr}
    \end{center}
\end{frame}

% _____________________________________________________________________________

\begin{frame}{Default execution space summary}
    \begin{center}
        \includegraphics[width=\textwidth]{default_execution_space.png}
    \end{center}
\end{frame}

% _____________________________________________________________________________

\begin{frame}[fragile]{Relation between execution space and memory space}
    \begin{itemize}
        \item That relation makes sense (both on same hardware)
        \item \texttt{memory\_space} attribute to access the memory space associated with an execution space
        \item \texttt{HostSpace} for memory space of the default host execution space
    \end{itemize}
    \pause
    \begin{center}
        % allow the table to overflow horizontally
        \makebox[\linewidth][c]{%
            \begin{tblr}[theme=kokkostable]{colspec=lll, row{2}={bg=lightmain}}
                & \SetCell[c=2]{l} Kokkos compiled on \\
                Memory space & CPU backend & GPU backend \\
                \texttt{Kokkos::DefaultExecutionSpace::memory\_space} & On CPU & On GPU \\
                \texttt{Kokkos::DefaultHostExecutionSpace::memory\_space} & On CPU & On CPU \\
            \end{tblr}%
        }
    \end{center}
\end{frame}

% _____________________________________________________________________________

\begin{frame}[fragile]{Anatomy of a nested loops}
    \begin{columns}
        \begin{column}{0.5\linewidth}
            \begin{minted}{C++}
                for (int i = 0; i < Nx; i++)
                    for (int j = 0; j < Ny; j++) {
                        A(i, j) = i + j;
                    }
            \end{minted}
        \end{column}
        \begin{column}{0.5\linewidth}
            \begin{itemize}
                \item 2 loop patterns
                \begin{itemize}
                    \item No extra code before/after the inner loop
                    \item Loops are collapsible (OpenMP \texttt{collapse})
                \end{itemize}
                \item 2 execution policies
                \begin{itemize}
                    \item No dependencies between them
                \end{itemize}
                \item 1 kernel
                \begin{itemize}
                    \item No iteration order dependencies
                    \item No side effects
                \end{itemize}
            \end{itemize}
        \end{column}
    \end{columns}
\end{frame}

% _____________________________________________________________________________

\begin{frame}[fragile]{How to have multidimensional loop policies?}
    \begin{columns}
        \begin{column}{0.66\linewidth}
            \begin{minted}{C++}
                Kokkos::MDRangePolicy<
                    Kokkos::Rank<N>
                > policy(
                    {start_index_0, /* ... */, start_index_n},
                    {end_index_0, /* ... */, end_index_n}
                );
            \end{minted}
        \end{column}
        \begin{column}{0.34\linewidth}
            \begin{itemize}
                \item \mintinline{C++}|MDRangePolicy| for nested loops (with \emph{MD} for multidimensional)
                \item \mintinline{C++}|Rank| template parameter to specify the dimension
                \item Max rank is 6
                \item \mintinline{C++}|start_index_*| and \mintinline{C++}|end_index_*| are the beginning (inclusive) and the end (exclusive) of each loop
            \end{itemize}
        \end{column}
    \end{columns}
\end{frame}

% _____________________________________________________________________________

\begin{frame}[fragile]{Anatomy of a Kokkos nested loop}
    \begin{columns}
        \begin{column}{0.6\linewidth}
            \begin{minted}{C++}
                Kokkos::parallel_for(
                    "my_loop",
                    Kokkos::MDRangePolicy<
                        Kokkos::Rank<2>
                    >(
                        {0, 0},
                        {Nx, Ny}
                    ),
                    KOKKOS_LAMBDA (int i, int j) {
                        A(i, j) = i + j;
                    }
                );
            \end{minted}
        \end{column}
        \begin{column}{0.4\linewidth}
            \begin{itemize}
                \item Still 1 \texttt{parallel\_for}
                \item No shortcut version!
                \item Kernel has 2 arguments
            \end{itemize}
        \end{column}
    \end{columns}
\end{frame}

\begin{frame}{Iteration order}
    \begin{center}
        \includegraphics[width=0.75\textwidth]{iterate_right_left.png}
    \end{center}
    \begin{itemize}
        \item Similarly with layout for views, iteration order is the way multidimensional indices are traversed by threads
        \item The default iterate order depends on the backend
        \begin{itemize}
            \item Order right on CPU
            \item Order left on GPU
        \end{itemize}
    \end{itemize}
\end{frame}

\begin{frame}[fragile]{Iteration order(s) for Kokkos}
    \begin{columns}
        \begin{column}{0.5\linewidth}
            \begin{center}
                \includegraphics[width=0.75\linewidth]{iterate_kokkos}
            \end{center}
        \end{column}
        \begin{column}{0.5\linewidth}
            \begin{itemize}
                \item Kokkos iterates by tiles, so there are 2 iteration orders
                \begin{description}[Outer order]
                    \item[Outer order] Order of the tiles within the iteration space
                    \item[Inner order] Order of the elements within the tile
                \end{description}
                \item The idea is to map with
                \begin{itemize}
                    \item GPU blocks and threads
                    \item CPU cache blocking
                \end{itemize}
                \item For best performance, the inner iteration order must match the memory layout of the accessed views
            \end{itemize}
        \end{column}
    \end{columns}
\end{frame}

\begin{frame}[fragile]{Iteration order(s) for Kokkos (API)}
    \begin{columns}
        \begin{column}{0.5\linewidth}
            \begin{minted}{C++}
                using Rank = Kokkos::Rank<
                    N,
                    // outer
                    Kokkos::Iterate::Default,
                    // inner
                    Kokkos::Iterate::Default
                >;
            \end{minted}
        \end{column}
        \begin{column}{0.5\linewidth}
            \begin{itemize}
                \item \texttt{Rank} to specify the iteration orders
                \item Possible values
                \begin{description}[\texttt{Iterate::Default}]
                    \item[\texttt{Iterate::Default}] Depends on the backend
                    \item[\texttt{Iterate::Left}] Left order
                    \item[\texttt{Iterate::Right}] Right order
                \end{description}
                \item (Again) for best performance, the inner iteration order must match the memory layout of the accessed views
            \end{itemize}
        \end{column}
    \end{columns}
\end{frame}

% _____________________________________________________________________________

\begin{exerciseframe}{Exercise 5: Multidimensional parallel loop}
    \begin{columns}
        \begin{column}{0.5\linewidth}
            \begin{center}
                \includegraphics[width=0.9\textwidth]{sleeping_otter.png}
            \end{center}

            Go to the exercise \githublink{\href{https://github.com/PaulGannay/cexa-kokkos-tutorials/tree/main/exercises/05_parallel_md_loop}{05\_parallel\_md\_loop}} and follow the instructions
        \end{column}
        \begin{column}{0.5\linewidth}
            \begin{block}{Goal of this exercise}
                \begin{itemize}
                    \item Update to a multidimensional policy
                    \item Ensure correctness of the results
                    \item Measure the performance for CPU and GPU
                    \item Experiment with different iteration orders
                \end{itemize}
            \end{block}
        \end{column}
    \end{columns}
\end{exerciseframe}

% _____________________________________________________________________________

\section{Parallel reduction}

% _____________________________________________________________________________

\begin{frame}[fragile]{Anatomy of a reduction loop}
    \begin{columns}
        \begin{column}{0.5\linewidth}
            \begin{minted}[texcomments]{C++}
                double result = 0;

                for (int i = 0; i < N; i++) {
                    result += A(i);
                }
            \end{minted}
        \end{column}
        \begin{column}{0.5\linewidth}
            \begin{itemize}
                \item One scalar variable collects something for each iteration (sum, max, min, etc.), to produce one final global result
                \pause
                \item Horrible truth: \mintinline{c++}{result} could \highlight{depend} on iteration order!
                \item This is due to floating point arithmetic not being associative and to unknown iteration order
            \end{itemize}
        \end{column}
    \end{columns}
\end{frame}

% _____________________________________________________________________________

\begin{frame}[fragile]{Anatomy of a Kokkos parallel reduction loop}
    \begin{columns}
        \begin{column}{0.5\linewidth}
            \begin{minted}{C++}
                // no initialization needed,
                // handled by parallel_reduce
                double result;

                Kokkos::parallel_reduce(
                    "my_reduction",
                    N,
                    KOKKOS_LAMBDA (
                        const int i,
                        double& local_result
                    ) {
                        local_result += i;
                    },
                    Kokkos::Sum<double>(result)
                );
            \end{minted}
        \end{column}
        \begin{column}{0.5\linewidth}
            \begin{itemize}
                \item \mintinline{c++}{parallel_reduce} to perform reductions
                \item Works like \texttt{parallel\_for} (e.g. execution policy, etc.)
                \item Reducer object (here \mintinline{c++}{Sum}), templated with the result type
                \item Extra last argument to the kernel: a reference to the local result
                \item Update local result in a manner consistent with global result
                \item The local result is write-only (it's a partial result)
            \end{itemize}
        \end{column}
    \end{columns}
\end{frame}

% _____________________________________________________________________________

\begin{frame}{About reducers}
    \begin{itemize}
        \item Kokkos provides implementations of Reducer classes for the most common operations:
        \begin{itemize}
            \item \texttt{Kokkos::Sum}: Sum of the elements
            \item \texttt{Kokkos::Max}: and \texttt{Kokkos::Min}: Maximum and minimum element
            \item \texttt{Kokkos::Prod}: Product of the elements
            \item \texttt{Kokkos::MaxLoc} and \texttt{Kokkos::MinLoc}: Maximum and minimum element with their index
            \item More \doclink{\href{https://kokkos.org/kokkos-core-wiki/ProgrammingGuide/Custom-Reductions-Built-In-Reducers.html}{in the doc}}
        \end{itemize}
        \item You can also create your own Reducer class, not covered in this talk
    \end{itemize}
\end{frame}

% _____________________________________________________________________________

\begin{frame}[fragile]{Sum reduction example}
    \begin{minted}{C++}
        double result;
        Kokkos::View<double*> A;
        ...
        Kokkos::parallel_reduce(
            "my_reduction",
            N,
            KOKKOS_LAMBDA (const int i, double& local_result) {
                local_result += A(i);
            },
            result // Only in the case of Sum, no need of the Reducer object
        );
    \end{minted}
\end{frame}

% _____________________________________________________________________________

\begin{frame}[fragile]{Max reduction example}
    \begin{minted}{C++}
        double max_value;
        Kokkos::View<double*> A;
        ...
        Kokkos::parallel_reduce(
            "my_reduction",
            N,
            KOKKOS_LAMBDA (const int i, double& local_max_value) {
                local_max_value = Kokkos::max(A(i), local_max_value);
                // local reduction consistent with the Reducer object
            },
            Kokkos::Max<double>(max_value)
        );
    \end{minted}

    \pause

    \structure{Note:} Notice that mathematical operators are in the \texttt{Kokkos} namespace

    \pause

    \structure{Beware:} \texttt{Kokkos::max} is a mathematical operator, \texttt{Kokkos::Max<double>} defines a Reducer object
\end{frame}

% _____________________________________________________________________________

\begin{frame}[fragile]{Min and index reduction example}
    \begin{columns}
        \begin{column}{0.55\linewidth}
            \begin{minted}[fontsize=\scriptsize]{C++}
                    using minloc_type =
                        Kokkos::MinLoc<double, int>::value_type;

                    minloc_type minloc;
                    Kokkos::parallel_reduce(
                        "min_loc_reduction",
                        N,
                        KOKKOS_LAMBDA (
                            int i,
                            minloc_type& local_minloc_value
                        ) {
                            if (A(i) < local_minloc_value.val) {
                                local_minloc_value.val = A(i);
                                local_minloc_value.loc = i;
                            }
                        },
                        Kokkos::MinLoc<double, int>(minloc)
                    );
                \end{minted}
        \end{column}
        \begin{column}{0.45\linewidth}
            \begin{itemize}
                \item Pay attention to the type of the intermediate and the global results
                \item Here, \texttt{Kokkos::MinLoc<double, int>::value\_type}: contains the value and the location
                \item In general, the type can be obtained as \texttt{value\_type} of the Reducer object
            \end{itemize}
        \end{column}
    \end{columns}
\end{frame}

% _____________________________________________________________________________

\begin{frame}[fragile]{Combining multiple reductions}
    \begin{itemize}
        \item Multiple Reducer objects can be specified in the same \texttt{parallel\_reduce}
    \end{itemize}
    \begin{minted}{C++}
        double sum_val,  min_val;

        Kokkos::parallel_reduce(
            "SumMinReduce", N,
            KOKKOS_LAMBDA (const int i, double& local_sum, double& local_min) {
                local_sum += A(i);
                local_min = Kokkos::min(A(i), local_min);
            },
            sum_val, Kokkos::Min<double>(min_val)
        );
    \end{minted}
\end{frame}

% _____________________________________________________________________________

\begin{exerciseframe}{Exercise 5: Parallel Reduction}
    \begin{columns}
        \begin{column}{0.5\linewidth}
            \begin{center}
                \includegraphics[width=0.9\textwidth]{sleeping_otter.png}
            \end{center}

            Go to the exercise \githublink{\href{https://github.com/CExA-project/cexa-kokkos-tutorials/tree/main/exercises/05_parallel_reduce}{06\_parallel\_reduce}} and follow the instructions
        \end{column}
        \begin{column}{0.5\linewidth}
            \begin{block}{Goal of this exercise}
                \begin{itemize}
                    \item Perform a parallel reduction
                    \item Measure the difference of performance between CPU and GPU
                    \item Increase the size of the problem to make the differences more visible
                    \item At which size the GPU is more efficient?
                \end{itemize}
            \end{block}
        \end{column}
    \end{columns}
\end{exerciseframe}

% _____________________________________________________________________________

\section{Parallel scans}

\begin{frame}[fragile]{Anatomy of a scan loop}
    \begin{columns}
        \begin{column}{0.55\linewidth}
            \begin{minted}[texcomments]{C++}
                result(0) = A(0);
                for (int i = 1; i < N; i++) {
                    result(i) = result(i - 1) + A(i);
                }
            \end{minted}
        \end{column}
        \begin{column}{0.45\linewidth}
            \begin{itemize}
                \item Cumulative reduction, similar to \texttt{numpy.cumsum}
                \item Inclusive (shown here) and exclusive scan algorithms
                \item Useful cumulative integrals and CSR matrix assembly
            \end{itemize}
        \end{column}
    \end{columns}
\end{frame}

\begin{frame}[fragile]{Anatomy of a Kokkos parallel scan loop}
    \begin{columns}
        \begin{column}{0.5\linewidth}
            \begin{minted}{C++}
                Kokkos::parallel_scan(
                    "my_scan",
                    N,
                    KOKKOS_LAMBDA(
                        int i,
                        int& update,
                        bool is_final
                    ) {
                        update += A(i);

                        if (is_final) {
                            result(i) = update;
                        }
                    }
                );
            \end{minted}
        \end{column}
        \begin{column}{0.5\linewidth}
            \begin{itemize}
                \item \mintinline{c++}{parallel_scan} to perform prefix reductions
                \item Works only with 1D execution policies
                \item In custom scan operations the functor needs to provide \texttt{init}, \texttt{join} and \texttt{operator()}
                \item Beware, the functor may be called multiple times (\texttt{is\_final})!
            \end{itemize}
        \end{column}
    \end{columns}
\end{frame}

\section{Function annotation}

% _____________________________________________________________________________

\begin{frame}[fragile]{Why do we need \texttt{KOKKOS\_LAMBDA}?}
    \begin{columns}
        \begin{column}{0.4\linewidth}
            \begin{minted}{C++}
                Kokkos::parallel_for(
                    "my_loop", N,
                    [=] (int i) {
                        A(i) = B(i) + C(i);
                    }
                );
            \end{minted}
            \begin{uncoverenv}<3->
                \structure{Question:} Why do we need to use \texttt{KOKKOS\_LAMBDA} inside parallel constructs rather than the standard C++ lambda syntax?
            \end{uncoverenv}
        \end{column}
        \begin{column}{0.6\linewidth}
            \begin{itemize}
                \item Let's create a \texttt{parallel\_for} without \texttt{KOKKOS\_LAMBDA}
                \item This code works when compiled with a CPU backend
                \pause
                \item But fails when compiled on CUDA backend with a cryptic error message:
                \begin{minted}[fontsize=\scriptsize]{text}
                    error: The closure type for a lambda ("lambda [](int)->void") cannot be used in the template argument type of a __global__ function template instantiation, unless the lambda is defined within a __device__ or __global__ function
                \end{minted}
            \end{itemize}
        \end{column}
    \end{columns}
\end{frame}

% _____________________________________________________________________________

\begin{frame}[fragile]{Function annotations}
    \begin{columns}
        \begin{column}{0.47\linewidth}
            \begin{minted}{cuda}
                __device__
                void device_function() {}

                __host__
                void host_function() {}

                __host__ __device__
                void host_device_function() {}

                auto lambda =
                  __host__ __device__ [] () {};
            \end{minted}
        \end{column}
        \begin{column}{0.55\linewidth}
            \begin{itemize}
                \item Some GPU backends (CUDA and HIP) need special annotations to know which code will be run on the host and which code will run on the device
                \item Lambdas also need to be annotated
                \item Function annotated with both \texttt{\_\_host\_\_} and \texttt{\_\_device\_\_} will be callable from the host and the device
                \item \texttt{\_\_host\_\_} is implicit if no annotation is given
                \item Neither \texttt{\_\_host\_\_} nor \texttt{\_\_device\_\_} are known by CPU backends, and will produce compilation errors
            \end{itemize}
        \end{column}
    \end{columns}
\end{frame}

% _____________________________________________________________________________

\begin{frame}[fragile]{Annotated functions calls}
    Calling a function with the wrong annotation will generate errors and warnings:
    \begin{center}
        \begin{tblr}[theme=kokkostable]{colspec=llll}
            \diagbox[linecolor=white,linewidth=2pt]{Caller}{Calling}& \texttt{\_\_host\_\_}  & \texttt{\_\_device\_\_} & both\\
            \texttt{\_\_host\_\_} & Ok & Error & Ok \\
            \texttt{\_\_device\_\_} & Error & Ok & Ok \\
          both & Warning & Error & Ok \\
        \end{tblr}
    \end{center}
\end{frame}

% _____________________________________________________________________________

\begin{frame}[fragile]{Kokkos annotation macros}
    \begin{columns}
        \begin{column}{0.45\linewidth}
            \begin{minted}{C++}
                KOKKOS_FUNCTION
                void device_function(){}

                KOKKOS_FUNCTION
                void host_function() {}

                KOKKOS_FUNCTION
                void host_device_function() {}

                auto lambda =
                    KOKKOS_LAMBDA () {};
            \end{minted}
        \end{column}
        \begin{column}{0.55\linewidth}
            \begin{itemize}
                \item \texttt{KOKKOS\_FUNCTION} and \texttt{KOKKOS\_LAMBDA} will have the following annotations:
                \begin{itemize}
                    \item \texttt{\_\_host\_\_ \_\_device\_\_} when compiling for a GPU backend
                    \item Nothing when compiling for a CPU backend
                \end{itemize}
            \end{itemize}
        \end{column}
    \end{columns}
\end{frame}

% _____________________________________________________________________________

\begin{frame}[fragile]{Annotated functions calls (with Kokkos macros)}
    \begin{center}
        \begin{tblr}[theme=kokkostable]{colspec=llll}
            \diagbox[linecolor=white,linewidth=2pt]{Caller}{Calling}& No macro  & \texttt{KOKKOS\_FUNCTION} & \texttt{KOKKOS\_LAMBDA} \\
            No macro & Ok & Ok & Ok \\
            \texttt{KOKKOS\_FUNCTION} & Error & Ok & Ok \\
            \texttt{KOKKOS\_LAMBDA} & Error & Ok & Ok \\
        \end{tblr}
    \end{center}
\end{frame}

\begin{frame}[fragile]{Use of annotated functions in Kokkos}
    \begin{columns}
        \begin{column}{0.5\linewidth}
            \begin{minted}{C++}
                KOKKOS_FUNCTION
                int function(int i) {
                    // ...
                }

                void compute() {
                    Kokkos::parallel_for(
                        "compute",
                        N,
                        KOKKOS_LAMBDA (int i) {
                            int a = function(i);
                            // ...
                        }
                    );
                }
            \end{minted}
        \end{column}
        \begin{column}{0.5\linewidth}
            \begin{itemize}
                \item You already know where to use \texttt{KOKKOS\_LAMBDA}
                \pause
                \item \texttt{KOKKOS\_FUNCTION} can be used to create free functions called in a lambda
            \end{itemize}
        \end{column}
    \end{columns}
\end{frame}

\begin{frame}[fragile]{From CUDA/HIP to Kokkos}
    \begin{columns}[T]
        \begin{column}{0.55\linewidth}
            \begin{center}
                \textbf{CUDA code}
            \end{center}

            \vspace{-1em}

            \begin{minted}[fontsize=\scriptsize]{cuda}
                auto lambda =
                    __host__ __device__ [=] (int i) {
                        // ...
                    };

                __global__
                void entrypoint() {
                    int i =
                        blockIdx.x * blockDim.x +
                        threadIdx.x;

                    if (i < N) lambda(i);
                }

                entrypoint<<<N / n_threads, n_threads>>>();
            \end{minted}
        \end{column}
        \begin{column}{0.45\linewidth}
            \begin{center}
                \textbf{Equivalent Kokkos code}
            \end{center}

            \vspace{-1em}

            \begin{minted}[fontsize=\scriptsize]{C++}
                auto lambda =
                    KOKKOS_LAMBDA (int i) {
                        // ...
                    };

                Kokkos::parallel_for(
                    "kernel",
                    N,
                    lambda
                );
            \end{minted}

            \pause

            \begin{itemize}
                \item Kokkos takes care of the entrypoint
                \item So you focus on the kernel
            \end{itemize}
        \end{column}
    \end{columns}
\end{frame}

% _____________________________________________________________________________

\section[Conclusion]{Conclusion}

% _____________________________________________________________________________

\begin{frame}{Conclusion}
    Congratulation, you have learned the basics of Kokkos!

    \begin{itemize}
        \item How to compile Kokkos
        \item How to create and manage a simple View
        \item How to use mirror Views and deep copy
        \item How to write a parallel loop
        \item How to extend the loop policy
        \item How to use nested parallel loops
        \item How to write a parallel reduction
        \item How to write a parallel scan
        \item How to use Kokkos macros for functions and lambdas
    \end{itemize}
\end{frame}

\end{document}
